رویدادها را به صورت زیر در نظر می‌گیریم:
\[
A:\ \text{ابتلا به سرطان ریه},\qquad
B_1:\ \text{تنفس طولانیِ هوای آلوده},\qquad
B_2:\ \text{بخارهای شیمیایی},\qquad
B_3:\ \text{سیگار کشیدن}
\]
با توجه به اینکه $\{B_1,B_2,B_3\}$ یک افراز است، از قانون احتمال کل داریم:
\begin{align*}
P(A) &= P(A|B_1)P(B_1)+P(A|B_2)P(B_2)+P(A|B_3)P(B_3)\\
&= (0.3)(0.7)+(0.5)(0.05)+(0.4)(0.25)\\
&= 0.21+0.025+0.10\\
&= 0.335
\end{align*}

حالا با قضیه‌ی \lr{Bayes}:
\begin{align*}
P(B_1|A) &= \frac{P(A|B_1)P(B_1)}{P(A)}
= \frac{(0.3)(0.7)}{0.335}
= \frac{0.21}{0.335}
= \frac{42}{67}
\approx 0.627
\end{align*}